\section*{Discussion}

Our findings suggest that persona instructions may change
how a model responds to safe and unsafe requests, but they
do not identify a consistently better persona. The neutral
baseline and children's tutor rarely refused safe requests,
yet more often complied with unsafe ones. Some professional
personas showed the opposite tendency. This illustrates why
refusing more often does not necessarily mean better safety
calibration. Useful model should distinguish between
requests that should be answered and the ones that should be
refused. Although the primary analysis detected overall
persona differences, no individual role differed significantly
from the neutral baseline after correction, and conclusions
depended on annotation choices.

A central finding concerns the difficulty of measuring
refusal behavior. The revised judge performed well on the
validation sample, but further review still identified
mistakes among the responses that directly determine
over-refusal and under-refusal rates. Persona-specific
language, such as fictional framing or professional
disclaimers, can make an answer appear more restrictive
than it actually is. Consequently, an apparent persona
effect may partly reflect how the evaluator interprets
that persona's style. The review procedure helped expose
this problem, but a small validation sample was insufficient
to establish reliability across all response types.

Separating deflection from refusal also changed the results. A response that says a request is outside
its professional scope provides no help, but does not
necessarily express a safety judgment. Counting such
responses as safety refusals would mix role adoption
with safety behavior. Conversely, deflecting an unsafe
request does not demonstrate that the model recognized
its harmfulness. Reporting deflection separately therefore
provides a clearer account of the model's behavior.

Several limitations restrict the conclusions. Later review
labels came from Claude rather than independent human
annotators, so corrections may reflect another model's
interpretation of ambiguous responses. Review also focused
on judge-identified potential errors; mistakes could remain
elsewhere despite agreement on the additional comparison
sample. The experiment used one quantized model, one
benchmark, and one greedy response per condition.
The 256-token limit may also have cut off content relevant
to annotation.

Future work should fix annotation rules before evaluation
and obtain independent human labels across personas and
response categories, including apparently correct judge
decisions. Testing additional models, alternative persona
wordings, and longer responses would help establish whether
the observed patterns generalize or new patterns emerge.
